% Options for packages loaded elsewhere
\PassOptionsToPackage{unicode}{hyperref}
\PassOptionsToPackage{hyphens}{url}
\PassOptionsToPackage{dvipsnames,svgnames,x11names}{xcolor}
\documentclass[
  10pt,
  fontset=fandol]{ctexart}
\usepackage{xcolor}
\usepackage[margin=16mm]{geometry}
\usepackage{amsmath,amssymb}
\setcounter{secnumdepth}{-\maxdimen} % remove section numbering
\usepackage{iftex}
\ifPDFTeX
  \usepackage[T1]{fontenc}
  \usepackage[utf8]{inputenc}
  \usepackage{textcomp} % provide euro and other symbols
\else % if luatex or xetex
  \usepackage{unicode-math} % this also loads fontspec
  \defaultfontfeatures{Scale=MatchLowercase}
  \defaultfontfeatures[\rmfamily]{Ligatures=TeX,Scale=1}
\fi
\usepackage{lmodern}
\ifPDFTeX\else
  % xetex/luatex font selection
\fi
% Use upquote if available, for straight quotes in verbatim environments
\IfFileExists{upquote.sty}{\usepackage{upquote}}{}
\IfFileExists{microtype.sty}{% use microtype if available
  \usepackage[]{microtype}
  \UseMicrotypeSet[protrusion]{basicmath} % disable protrusion for tt fonts
}{}
\makeatletter
\@ifundefined{KOMAClassName}{% if non-KOMA class
  \IfFileExists{parskip.sty}{%
    \usepackage{parskip}
  }{% else
    \setlength{\parindent}{0pt}
    \setlength{\parskip}{6pt plus 2pt minus 1pt}}
}{% if KOMA class
  \KOMAoptions{parskip=half}}
\makeatother
\setlength{\emergencystretch}{3em} % prevent overfull lines
\providecommand{\tightlist}{%
  \setlength{\itemsep}{0pt}\setlength{\parskip}{0pt}}
\usepackage{amsmath,amssymb,mathtools,needspace}
\xeCJKDeclareCharClass{CJK}{"2032,"0400->"04FF,"2160->"216B,"2460->"2473,"25A0,"25C7,"25CB,"25CF}
\IfFontExistsTF{Arial Unicode MS}{\xeCJKsetup{AutoFallBack=true}\setCJKfallbackfamilyfont{\CJKrmdefault}{Arial Unicode MS}}{}
\setcounter{secnumdepth}{0}
\setlength{\emergencystretch}{2em}
\allowdisplaybreaks
\usepackage{bookmark}
\IfFileExists{xurl.sty}{\usepackage{xurl}}{} % add URL line breaks if available
\urlstyle{same}
\hypersetup{
  pdftitle={代数方程求根},
  colorlinks=true,
  linkcolor={Maroon},
  filecolor={Maroon},
  citecolor={Blue},
  urlcolor={Blue},
  pdfcreator={LaTeX via pandoc}}

\title{代数方程求根}
\author{}
\date{}

\begin{document}
\maketitle

\subsection{6.5 代数方程求根}\label{ux4ee3ux6570ux65b9ux7a0bux6c42ux6839}

如果 \(f(x)\) 是多项式，则 \(f(x)=0\) 特别地称为代数方程。前面介绍的求根方法原则上也适用于解代数方程，但由于多项式的特殊性，可以针对其特点提供更为有效的算法。

\subsubsection{6.5.1 多项式求值的秦九韶算法}\label{ux591aux9879ux5f0fux6c42ux503cux7684ux79e6ux4e5dux97f6ux7b97ux6cd5}

多项式的一个重要特点是求值方便。

设给定多项式

\[
f(x)=a_0x^n+a_1x^{n-1}+\cdots+a_{n-1}x+a_n,
\]

其中系数 \(a_i\ (0\leqslant i\leqslant n)\) 均为实数。用一次式 \(x-x_0\) 除 \(f(x)\)，商记作 \(P(x)\)，余数显然等于 \(f(x_0)\)，即有

\[
f(x)=f(x_0)+(x-x_0)P(x).
\tag{6.5.1}
\]

现在具体确定 \(P(x)\) 与 \(f(x_0)\)。令

\[
p(x)=b_0x^{n-1}+b_1x^{n-2}+\cdots+b_{n-2}x+b_{n-1},
\]

代入式（6.5.1），并比较两端同次幂的系数，得

\[
\begin{cases}
a_0=b_0,\\
a_i=b_i-x_0b_{i-1},\quad1\leqslant i\leqslant n-1,\\
a_n=f(x_0)-x_0b_{n-1}.
\end{cases}
\]

从而有

\[
\begin{cases}
b_0=a_0,\\
b_i=a_i+x_0b_{i-1},\quad1\leqslant i\leqslant n,\\
f(x_0)=b_n.
\end{cases}
\tag{6.5.2}
\]

\href{../../source-images/pdf-172.jpeg}{原页}

这里提供的一种计算函数值 \(f(x_0)\) 的有效算法称为秦九韶法①。这种算法的优点是计算量小，结构紧凑，容易编制计算程序。

进一步考察 \(f(x)\) 的 Taylor 展开式

\[
f(x)=f(x_0)+f'(x_0)(x-x_0)+\frac{f''(x_0)}{2!}(x-x_0)^2+\cdots+\frac{f^{(n)}(x_0)}{n!}(x-x_0)^n.
\]

将它表示为式（6.5.1）的形式，则式中

\[
P(x)=f'(x_0)+\frac{f''(x_0)}{2!}(x-x_0)+\cdots+\frac{f^{(n)}(x_0)}{n!}(x-x_0)^{n-1}.
\]

由此可见，导数 \(f'(x_0)\) 又可看作 \(P(x)\) 用因式 \(x-x_0\) 相除得出的余数，即有

\[
P(x)=f'(x_0)+(x-x_0)Q(x),
\]

其中 \(Q(x)\) 是 \(n-2\) 次多项式。设

\[
Q(x)=c_0x^{n-2}+c_1x^{n-3}+\cdots+c_{n-3}x+c_{n-2},
\]

那么，再用秦九韶算法又可求出值 \(f'(x_0)\)，对应于式（6.5.2），这里计算公式是

\[
\begin{cases}
c_0=b_0,\\
c_i=b_i+x_0c_{i-1},\quad1\leqslant i\leqslant n-1,\\
f'(x_0)=c_{n-1}.
\end{cases}
\tag{6.5.3}
\]

继续这一过程，可以依次求出 \(f(x)\) 在点 \(x_0\) 的各阶导数。

\subsubsection{6.5.2 代数方程的 Newton 法}\label{ux4ee3ux6570ux65b9ux7a0bux7684-newton-ux6cd5}

再就多项式方程

\[
f(x)=a_0x^n+a_1x^{n-1}+\cdots+a_{n-1}x+a_n=0
\]

考察 Newton 公式

\[
x_{k+1}=x_k-\frac{f(x_k)}{f'(x_k)}.
\tag{6.5.4}
\]

据式（6.5.2）与式（6.5.3），式（6.5.4）中的函数值 \(f(x_k)\) 和导数值 \(f'(x_k)\) 均可方便地求出

\[
\begin{cases}
b_0=a_0,\\
b_i=a_i+x_kb_{i-1},\quad1\leqslant i\leqslant n,\\
f(x_k)=b_n,
\end{cases}
\tag{6.5.5}
\]

\[
\begin{cases}
c_0=b_0,\\
c_i=b_i+x_kc_{i-1},\quad1\leqslant i\leqslant n-1,\\
f'(x_k)=c_{n-1}.
\end{cases}
\tag{6.5.6}
\]

① 这种算法是我国宋代的一位数学家秦九韶最先提出的。外国文献中通常称这种算法为 Horner 算法，其实 Horner 的工作比秦九韶晚了五六个世纪。

\href{../../source-images/pdf-173.jpeg}{原页}

\subsubsection{6.5.3 劈因子法}\label{ux5288ux56e0ux5b50ux6cd5}

如果能从多项式 \(f(x)=a_0x^n+a_1x^{n-1}+\cdots+a_{n-1}x+a_n\) 中分离出一个二次因式

\[
\omega^*(x)=x^2+u^*x+v^*,
\]

就能获得它的一对共轭复根。劈因子法的基本思想是，从某个近似的二次因子

\[
\omega(x)=x^2+ux+v
\]

出发，用某种迭代过程使之逐步精确化。

用二次式 \(\omega(x)\) 除 \(f(x)\)，商记作 \(p(x)\)，它是个 \(n-2\) 次多项式，余式为一次式，记作 \(r_0x+r_1\)，即

\[
f(x)=(x^2+ux+v)p(x)+r_0x+r_1.
\tag{6.5.7}
\]

显然，\(r_0,r_1\) 均为 \(u,v\) 的函数 \(\begin{cases}r_0=r_0(u,v),\\r_1=r_1(u,v).\end{cases}\) 劈因子法的目的就是逐步修改 \(u,v\) 的值，使余数 \(r_0,r_1\) 变得很小。考察方程

\[
\begin{cases}
r_0(u,v)=0,\\
r_1(u,v)=0.
\end{cases}
\tag{6.5.8}
\]

这是关于 \(u,v\) 的非线性方程组，设它有解 \((u^*,v^*)\)，将 \(r_0(u^*,v^*)=0,r_1(u^*,v^*)=0\) 的左端在 \((u,v)\) 展开到一阶项，则有

\[
\begin{cases}
r_0+\dfrac{\partial r_0}{\partial u}(u^*-u)+\dfrac{\partial r_0}{\partial v}(v^*-v)\approx0,\\
r_1+\dfrac{\partial r_1}{\partial u}(u^*-u)+\dfrac{\partial r_1}{\partial v}(v^*-v)\approx0.
\end{cases}
\]

这样，用 Newton 法的处理思想将非线性方程组（6.5.8）线性化，归结得到下列线性方程组：

\[
\begin{cases}
r_0+\dfrac{\partial r_0}{\partial u}\Delta u+\dfrac{\partial r_0}{\partial v}\Delta v=0,\\
r_1+\dfrac{\partial r_1}{\partial u}\Delta u+\dfrac{\partial r_1}{\partial v}\Delta v=0.
\end{cases}
\tag{6.5.9}
\]

从方程组（6.5.9）解出增量 \(\Delta u,\Delta v\)，即可得到改进的二次因式

\[
\omega(x)=x^2+(u+\Delta u)x+v+\Delta v.
\]

以下具体说明如何计算方程组（6.5.9）的各个系数。

（1）计算 \(r_0\) 和 \(r_1\)。将

\[
P(x)=b_0x^{n-2}+b_1x^{n-3}+\cdots+b_{n-3}x+b_{n-2}
\]

代入式（6.5.7），并比较各次幂的系数，易知

\[
\begin{cases}
a_0=b_0,\\
a_1=b_1+ub_0,\\
a_i=b_i+ub_{i-1}+vb_{i-2},\quad2\leqslant i\leqslant n-2,\\
a_{n-1}=ub_{n-2}+vb_{n-3}+r_0,\\
a_n=vb_{n-2}+r_1,
\end{cases}
\]

\href{../../source-images/pdf-174.jpeg}{原页}

于是 \(r_0,r_1\) 的计算公式为

\[
\begin{cases}
b_0=a_0,\\
b_1=a_1-ub_0,\\
b_i=a_i-ub_{i-1}-vb_{i-2},\quad2\leqslant i\leqslant n,\\
r_0=b_{n-1},\\
r_1=b_n+ub_{n-1}.
\end{cases}
\tag{6.5.10}
\]

（2）计算 \(\dfrac{\partial r_0}{\partial v},\dfrac{\partial r_1}{\partial v}\)。对式（6.5.7）关于 \(v\) 求导，得

\[
P(x)=-(x^2+ux+v)\frac{\partial P}{\partial v}+s_0x+s_1,
\tag{6.5.11}
\]

其中

\[
s_0=-\frac{\partial r_0}{\partial v},\quad s_1=-\frac{\partial r_1}{\partial v}.
\tag{6.5.12}
\]

可见，用 \(x^2+ux+v\) 除 \(P(x)\)，作为余式可以得到 \(s_0x+s_1\)。由于 \(P(x)\) 是 \(n-2\) 次多项式，这里商 \(\dfrac{\partial P}{\partial v}\) 是 \(n-4\) 次多项式，记

\[
\frac{\partial P}{\partial v}=c_0x^{n-4}+c_1x^{n-5}+\cdots+c_{n-5}x+c_{n-4},
\]

则相当于式（6.5.10），这里有

\[
\begin{cases}
c_0=b_0,\\
c_1=b_1-ub_0,\\
c_i=b_i-uc_{i-1}-vc_{i-2},\quad2\leqslant i\leqslant n-2,\\
s_0=c_{n-3},\\
s_1=c_{n-2}+uc_{n-3},
\end{cases}
\]

而按式（6.5.12）得

\[
\frac{\partial r_0}{\partial v}=-s_0,\quad\frac{\partial r_1}{\partial v}=-s_1.
\]

（3）计算 \(\dfrac{\partial r_0}{\partial u},\dfrac{\partial r_1}{\partial u}\)。对式（6.5.7）关于 \(u\) 求导，得

\[
xP(x)=-(x^2+ux+v)\frac{\partial P}{\partial u}-\frac{\partial r_0}{\partial u}x-\frac{\partial r_1}{\partial u}.
\]

另外，由式（6.5.11）有

\[
\begin{aligned}
xP(x)&=-(x^2+ux+v)x\frac{\partial P}{\partial v}+(s_0x+s_1)x\\
&=-(x^2+ux+v)\left(x\frac{\partial P}{\partial v}-s_0\right)-(us_0-s_1)x-vs_0.
\end{aligned}
\]

比较以上两式可知

\[
\frac{\partial r_0}{\partial u}=us_0-s_1,\quad\frac{\partial r_1}{\partial u}=vs_0.
\]

\begin{quote}
校注（agent 补充）：本页原文把用 \(x^2+ux+v\) 除 \(P(x)\) 的商记为 \(\partial P/\partial v\)，并给出以 \(c_0=b_0\) 开始的递推。式（6.5.11）本身表明该商应为 \(-\partial P/\partial v\)，因此商的叙述和紧随其后的展开式疑漏印负号。转录保留原文，后续的余数 \(s_0,s_1\) 及其递推亦按原页保留。
\end{quote}

\subsection{本节覆盖原页的校注记录}\label{ux672cux8282ux8986ux76d6ux539fux9875ux7684ux6821ux6ce8ux8bb0ux5f55}

下列记录按原页关联，含同页相邻小节的校注；计算时同时读取上文校注和完整原页。

\begin{itemize}
\tightlist
\item
  \texttt{ch06-E005}：原书 PDF 页 174
\end{itemize}

\end{document}
